%!TEX root = ../../main.tex
%% ===============================================================
%% 7.5 시공간 모델
%% 파일: chapters/07_models/05_spatiotemporal.tex
%% 상위: chapters/07_models/chapter.tex
%% 정의 라벨: sec:model-st
%% 참조 라벨: -
%% 포함 객체: -
%% 인용: yan2018spatial
%% 상태: 초안 (docs/OUTLINE.md 참고)
%% ===============================================================

\section{시공간 모델}
\label{sec:model-st}

\begin{itemize}[leftmargin=1.5em]
  \item \textbf{ST-GNN}: 구간마다 GNN을 적용한 뒤 시간 축을 GRU 또는 주의 풀링으로 요약.
  \item \textbf{ST-GCN} \cite{yan2018spatial}: 그래프 합성곱과 시간 합성곱 블록을 교대로 쌓음.
  \item \textbf{ST-Mamba}: GNN 공간 인코딩 후 Mamba 시간 인코딩.
  \item \textbf{다중 스케일 동적 그래프}: $A = \sum_k w_k A(\sigma_k)$의 학습 가중 커널 합.
  \item \textbf{EventXAttn}: 사건 토큰 임베딩 $e_i = \mathrm{Embed}(\mathrm{type}_i) + \mathrm{Linear}(\mathrm{cont}_i)$를 질의로, 그래프 인코딩된 플레이어 상태를 키·값으로 하는 교차 주의 후 중요도 가중 풀링.
\end{itemize}
